\section{Podsumowanie}

\subsection{Główne wnioski}

\begin{enumerate}
    \item \textbf{Przezwyciężenie ograniczeń miary Jordana}:
          miara Jordana, oparta na skończonych pokryciach, nie mierzy nawet tak prostych zbiorów jak $\Q \cap [0, 1]$. Przejście do przeliczalnych pokryć i $\sigma$-addytywności sprawia, że każdy zbiór przeliczalny ma miarę zero, a zbiory mierzalne tworzą $\sigma$-algebrę zawierającą $\BorelRn$.

    \item \textbf{Zupełność}:
          przestrzeń $(\Rn, \Leb{\Rn}, \lambda^n)$ jest zupełna, tzn.\ każdy podzbiór zbioru miary zero jest mierzalny; $\Leb{\Rn}$ jest uzupełnieniem $\BorelRn$.

    \item \textbf{Topologiczna natura mierzalności}:
          każdy zbiór mierzalny różni się od pewnego zbioru typu $G_\delta$ (oraz od pewnego zbioru typu $F_\sigma$) o zbiór miary zero, a jego miarę można przybliżać miarami zbiorów otwartych z zewnątrz i zwartych od wewnątrz.

    \item \textbf{Granice mierzalności}:
          zbiór Vitalego pokazuje, że przy aksjomacie wyboru nie istnieje nietrywialna, $\sigma$-addytywna i niezmiennicza na translacje miara określona na całym $\mathcal{P}(\Rn)$.
\end{enumerate}

\subsection{Znaczenie miary Lebesgue'a w analizie matematycznej}

Najważniejszą motywacją konstrukcji miary Lebesgue'a jest całka Lebesgue'a. W odróżnieniu od całki Riemanna, która dzieli dziedzinę funkcji na przedziały, całka Lebesgue'a dzieli zbiór wartości funkcji i mierzy przeciwobrazy otrzymanych przedziałów. Wymaga to, by te przeciwobrazy były zbiorami mierzalnymi, a im bogatsza jest klasa zbiorów mierzalnych, tym więcej funkcji można całkować. Przykładem jest funkcja Dirichleta, czyli funkcja charakterystyczna zbioru $\Q \cap [0,1]$: nie jest ona całkowalna w sensie Riemanna, natomiast jej całka Lebesgue'a istnieje i jest równa zeru, ponieważ zbiór liczb wymiernych ma miarę Lebesgue'a zero. Dzięki $\sigma$-addytywności miary całka Lebesgue'a spełnia też silne twierdzenia o przechodzeniu do granicy pod znakiem całki, których brakuje w teorii Riemanna.

Na tym fundamencie zbudowano dużą część współczesnej analizy: przestrzenie funkcji całkowalnych, w których zupełność miary przekłada się na zupełność przestrzeni, oraz analizę Fouriera, w której kluczową rolę odgrywa niezmienniczość miary Lebesgue'a względem translacji.

Teoria miary stała się również językiem rachunku prawdopodobieństwa. W 1933 roku A.~Kołmogorow zdefiniował prawdopodobieństwo jako miarę unormowaną na $\sigma$-algebrze zdarzeń, a miara Lebesgue'a na odcinku $[0,1]$ jest naturalnym modelem rozkładu jednostajnego.
